% Independent mathematical audits passed: root, path_stability, and random_hosts.
\section{A linear factor for two-cycle packings}
\label{sec:linear-two-cycle-packings}

This refinement improves the fourth-degree polynomial in Lemma~\ref{lem:high-path-two-cycle-packings} to a linear factor.
It uses the same marked layers and matching choices.
The improvement comes from summing the possible cycle intervals with their geometric weights.

\begin{lemma}[Two intervals with bounded overlap]
\label{lem:bounded-overlap-interval-sum}
Fix a nonnegative integer $B$.
There is a constant $c_B$ such that, for every integer $D\ge0$,
\[
 \sum_{\substack{I,J\text{ nonempty integer intervals in }[0,D]\\
                  |I\cap J|\le B}}
 2^{(|I|+|J|)/2}
 \le c_B(D+1)2^{D/2}.
\]
The intervals in this sum are ordered.
\end{lemma}

\begin{proof}
Put $u=\sqrt2$.
The one-interval sum satisfies
\[
 \sum_{0\le i\le j\le D}u^{j-i+1}
 \le\frac{u}{(1-u^{-1})^2}u^D<17u^D.
\]
Order the left endpoints, at a cost of at most two.
There are three cases.

First suppose that the intervals are disjoint, with the first before the second.
Let their outer endpoints be $i,l$, and let $g$ count the intervening unused positions.
Their total length is $l-i+1-g$.
For fixed $i,l,g$, there are at most $D+1$ choices of the split position.
Summing the geometric factor $u^{-g}$ gives a bound
\[
 \frac{17}{1-u^{-1}}(D+1)u^D.
\]

Next suppose that the intervals overlap without one containing the other.
Let $q\le B$ be their overlap length.
Their total length is the outer span plus $q$.
For fixed outer endpoints and $q$, there are at most $D+1$ meeting positions.
The resulting bound is
\[
 17\left(\sum_{q=1}^{B}u^q\right)(D+1)u^D.
\]
Finally, if one interval contains the other, the inner interval has length at most $B$.
For each outer interval and inner length, there are at most $D+1$ placements.
The same bound applies.
Adding the three cases proves the lemma.
\end{proof}

\begin{theorem}[Linear two-cycle packing bound]
\label{thm:linear-high-path-two-cycle-packings}
For every fixed $0<\varepsilon\le1/2$, there is a constant $B_\varepsilon$ such that the following holds.
Let $N$ be a normalized two-pole of rank $b$ with
\[
 P(N)\ge(1/2+\varepsilon)2^b.
\]
The number of nonempty unions of at most two vertex-disjoint cycles in $N$ is at most
\[
 B_\varepsilon(b+1)2^{b/2}.
\]
\end{theorem}

\begin{proof}
First suppose that $N$ is bridgeless.
Use the shortest terminal path and initial whole-layer marking from Lemma~\ref{lem:thin-layer-cycle-count}.
Additive stability bounds both $\Delta$ and $H$ by $C_\varepsilon$.
Thus the initial marked set has at most $10C_\varepsilon+2$ vertices and meets at most that many layers.
Write $D$ for the maximum distance layer, so $D\le b+2C_\varepsilon$.

Each cycle has a nonempty interval of distance layers.
Consider two vertex-disjoint cycles with intervals $I,J$.
Any layer strictly inside both intervals must contain at least two vertices from each cycle.
An initially ordinary layer has only two vertices, so it cannot have this property.
At most four layers are interval endpoints.
Consequently,
\[
 |I\cap J|\le \lceil10C_\varepsilon+2\rceil+4=:B.
\]

Fix the two intervals, and mark all their endpoint layers.
This adds at most eight vertices to the initial marked set.
Fix all selected edges incident with marked vertices, at a cost bounded by a constant depending only on $\varepsilon$.
Every remaining ordinary layer is unused or carries the two through strands of exactly one cycle.
The intervals determine which possibility applies.
Across consecutive occupied ordinary layers, the selected edges form a perfect matching.
Two-choice cuts cannot be consecutive, by the same cubic degree argument as before.
Their number is at most $(|I|+|J|)/2$.
The interval data, marked incidences, and matching choices determine the union of cycles.

Lemma~\ref{lem:bounded-overlap-interval-sum} now bounds two-cycle unions by
\[
 O_\varepsilon((D+1)2^{D/2})
 =O_\varepsilon((b+1)2^{b/2}).
\]
The one-cycle contribution is $O_\varepsilon(2^{b/2})$ by Theorem~\ref{thm:high-path-cycle-tradeoff}.
This proves the bridgeless case.

For a general normalized two-pole, use its series decomposition into bridgeless pieces.
Let their positive ranks be $b_1,\ldots,b_q$, with sum $b$.
Every piece retains the required normalized path density.
Packings contained in one piece contribute at most
\[
 B'_\varepsilon\sum_i(b_i+1)2^{b_i/2}
 \le B'_\varepsilon(b+2)2^{b/2}.
\]
Here we used $\sum_i2^{b_i/2}\le2\,2^{b/2}$.

Packings with one cycle in each of two pieces contribute at most
\[
 K_\varepsilon^2\sum_{i<j}2^{(b_i+b_j)/2}.
\]
For $q\ge2$,
\[
 \frac{\sum_{i<j}2^{(b_i+b_j)/2}}{2^{b/2}}
 \le\binom q2\,2^{-(q-2)/2}\le4.
\]
The last bound follows by checking $2\le q\le6$ and then taking successive ratios.
Increasing $B_\varepsilon$ absorbs both contributions.
\end{proof}

\begin{corollary}[Linear cross-sum forest bound]
\label{cor:linear-four-pole-cross-sum-forest}
For each fixed $\eta>0$, there is a constant $A_\eta$ such that every connected four-pole $K$ of rank $h$ satisfies
\[
 \sum_{i,j=1}^2P_{L_iR_j}(K)\ge(2+\eta)2^h
 \quad\Longrightarrow\quad
 F(K)\le A_\eta(h+3)2^{h/2}.
\]
Here $F(K)$ counts all vertex-disjoint path pairs connecting its four ports.
\end{corollary}

\begin{proof}
Cap the two sides by new terminals, as in Corollary~\ref{cor:four-pole-cross-sum-forest}.
The resulting two-pole has rank $h+2$ and normalized path density at least $1/2+\eta/4$.
Each four-terminal forest maps injectively to one cycle or two vertex-disjoint cycles in this two-pole.
Apply Theorem~\ref{thm:linear-high-path-two-cycle-packings}.
\end{proof}

The linear factor in the packing theorem is necessary.
Take a capped ladder consisting of $m$ square blocks.
Its terminal paths are saturated and its rank is $b=2m+1$.
Count two disjoint cycles, one through each cap terminal.
If the left cycle turns in square $i$ and the right cycle turns in square $j>i$, their number is $2^{m+i-j+1}$.
Summing these choices gives
\[
 \sum_{1\le i<j\le m}2^{m+i-j+1}
 =(m-2)2^{m+1}+4
 =\Theta(b2^{b/2}).
\]
The turns must occur in different squares because two same-square caps cannot be vertex-disjoint.
Thus a bound $O_\varepsilon(2^{b/2})$ without a rank factor would fail even at full terminal-path density.
